Let $A$ be a local ring, $m$ its maximal ideal, $q$ an ideal which is primary for $m$. If $A/m$ is an infinite field, there exists an ideal $q' \subset q$, generated by a system of parameters and such that $e(q') = e(q)$.

Proof. We proceed by induction on the dimension $d$ of $A$. For $d = 0$, every proper ideal $q$ of $A$ is nilpotent, whence all the characteristic polynomials $\bar{P}_q(n)$ are equal to the constant $\ell(A)$ ($= \text{length}$ of the $A$-module $A$); we may thus take $q' = (0)$, since we agree that (0) is generated by the empty set, which is thus the only system of parameters of $A$.

We now pass to the case $d = 1$. Note that, if $d = 1$, then $m$ is the

local ring $A/Ax$ has dimension 0. Consequently, for $n$ large, $\bar{P}_{\mathfrak{q}/Ax}(n)$ is the length of $A/Ax$. This is also the length of $Ax^{n-1}/Ax^n$ since $A/Ax$ and $Ax^{n-1}/Ax^n$ are isomorphic under the mapping $z + Ax \rightarrow zx^{n-1} + Ax^n$ (x not being a zero divisor). We therefore have $e(\mathfrak{q}) = e(Ax)$ in this case.

Still in case $d=1$, we now assume that $m$ is an imbedded prime ideal of (0) (i.e., that all the elements of $m$ are zero-divisors in $A$). Then the annihilator $a$ of $x$ is an ideal which, if considered as an $A$-module, has a finite length $s$ (since it is annihilated by $Ax$ and since, by Corollary 1 of Theorem 20, § 9, $A/Ax$ has dimension zero). We consider the factor ring $A^\star = A/\mathfrak{a}$. For every open ideal $\mathfrak{v}$ of $A$, we have $\mathfrak{v}^n \cap \mathfrak{a} = (0)$ for $n$ large since $\bigcap_{n=0}^{\infty} (\mathfrak{v}^n \cap \mathfrak{a}) = (0)$ and since $\mathfrak{a}$ is an $A$-module of finite length. Thus, if we set $\mathfrak{v}^\star = (\mathfrak{v} + \mathfrak{a})/\mathfrak{a}$, we have $P_{\mathfrak{v}^\star}(n) = \ell(A^\star/\mathfrak{v}^\star n) = \ell(A/\mathfrak{v}^n) - \ell((\mathfrak{a} + \mathfrak{v}^n)/\mathfrak{v}^n) = \ell(A/\mathfrak{v}^n) - \ell(\mathfrak

REMARK. If the ideal $q$ is generated by a system of parameters $\{y_1, \cdots, y_d\}$, it is also generated by a system of parameters $\{x_1, \cdots, x_d\}$ such that

$$e(q) = e(q/Ax_1) = \cdots = e(q/(Ax_1 + \cdots + Ax_{d-1})).$$

In fact, there exists, by a reasoning similar to the one given in Remark 2 to Lemma 5 (§ 8), an element $x_1$ of $q$ which is superficial of order 1 for $q$ and which may be written in the form $x_1 = uy_1 + a_2y_2 + \cdots + a_dy_d$ with $u \notin \mathfrak{m}$. (Observe that $[\mathfrak{m}y_1 + (y_2, \cdots, y_d) + q^2] / q^2$ is a proper submodule of $q / q^2$.) Then, by the last part of the proof of Theorem 22, we have $e(q) = e(q/Ax_1)$, and $q$ is also generated by $\{x_1, y_2, \cdots, y_d\}$. We operate in the same way in $A/Ax_1$, $A/(Ax_1 + Ax_2)$, etc., up to $x_{d-1}$. We take $x_d$ to be $y_d$. Now we consider the one-dimensional local ring $B = A/(Ax_1 + \cdots + Ax_{d-1})$ and the residue class $x^\star$ of $x_d$ in $B$. We have $e(q) = e(Bx^\star)$ by relation (3), and $x^\star$ is a parameter in $B$. If $x^\star$ is not a zero-divisor in $B$ (i.e., if the maximal ideal of $B$ is not entirely composed of zero divisors) then the modules $B/Bx^\star$ and $Bx^{\star n}/Bx^{\star n+1}$ are isomorphic (under the mapping $z + Bx

<n-q. These products generate an $(A/q)$-module whose length is at least $(n-q+d-1)$. Then the formula $\ell(A/q)^n = \ell(B/(\mathfrak{X}^n + \mathfrak{I}))$ implies that $P_q(n) \leq \ell(A/q)(n+d-1) - (n-q+d-1)$, whence $e(q) < \ell(A/q)$. This proves our second assertion. The converse is obvious.

Notice that, if $e(q) = \ell(A/q)$, the function $P_q(n)$ is a polynomial in $n$ for all values of $n$ (and not only for the large ones). Furthermore, since the initial form $X_i$ of $x_i$ is not a zero divisor in $G_q(A) = B$, we have $q^n:x_i = q^{n-1}$ for every $n$ and every $i$. As noticed in the remark following Theorem 22, this always happens if $A$ is a regular local ring.

We conclude this section with the proof of a theorem which not only can be used in certain cases for the computation of multiplicities, but also gives information on the behavior of multiplicities under finite integral extensions. This theorem is the algebraic counterpart of the projection formula for intersection cycles in Algebraic Geometry: